\section{Limitations and future work}
\label{sec:limits}

\fillin{Draft: extend with anything the experiments reveal.}

\paragraph{Audit.}
\begin{itemize}
  \item One person coded the papers, with AI assistance. Reliability is estimated only by a blind re-score of 7 papers.
  \item Eight core papers could be read only as abstracts.
  \item The search for earlier critical reviews was not completed.
\end{itemize}

\paragraph{Data.}
\begin{itemize}
  \item SMD has no anomaly types, so E2 and E3 give overall results only, and the shape-versus-volume question rests on injections. SMD is also open to the benchmark critiques in~\cite{Q04}. We use it because most TDA papers do, which lets us compare with them directly.
  \item CTU-13 has only 7 distinct normal hosts, 4 of which appear in every scenario. One scenario (10) lacks flow labels. E4 is therefore supporting evidence only.
\end{itemize}

\paragraph{Design.}
\begin{itemize}
  \item E5 is a mechanism test by design. Shape injections keep each channel's mean and variance, and volume injections keep the coupling between channels. It shows what each feature set \emph{can} detect, not how often such anomalies occur in practice.
  \item We use one detector with fixed defaults. Another detector might use topological features differently.
  \item F3 is one reasonable, untuned topology pipeline. A negative result for F3 is not a negative result for every possible topological method.
\end{itemize}

\paragraph{Future work.}
\begin{itemize}
  \item Test against an adversary who knows the features; no paper does (R9).
  \item Replicate I03 on DARPA OpTC with its original hypergraph construction.
  \item Run analyst studies for the exploratory uses (Mapper, motifs), which detection metrics do not capture.
  \item Measure streaming cost at realistic event rates, for Fluent Bit-style deployments.
\end{itemize}
